\documentclass[11pt]{article}
\usepackage[margin=1in]{geometry}
\usepackage[T1]{fontenc}
\usepackage{lmodern}
\usepackage{amsmath,amssymb}
\usepackage[hidelinks]{hyperref}

\title{Literature Status for arXiv:1501.01685:\\
The Riesz--Kantorovich Modulus Question}
\author{Run fa\_banach\_001, agent\_lane\_05}
\date{11 August 2026}

\begin{document}
\maketitle

\section*{Status}

This is a \texttt{literature\_already\_answered} packet.  Troitsky and
Xanthos recall a classical open problem in their paper \emph{Spaces of regular
abstract martingales}, arXiv:1501.01685 \cite{TX15}.  Michael Elliott later
answered that problem negatively in \emph{The Riesz--Kantorovich formulae}
\cite{Elliott19}.

\section*{Source Question}

In Section~2, on page~5 of the locally rendered source PDF, Troitsky and
Xanthos write that it is open whether the modulus of an operator is always
given by the Riesz--Kantorovich formula.  More explicitly, if a linear
operator $T:X\to Y$ between vector lattices has a modulus $|T|$ in the ordered
space of operators, must one have, for every $x\in X_+$,
\[
 |T|x=\sup T[-x,x]?
\]
The source cites Aliprantis--Tourky for this problem and then turns to the
analogous Krickeberg formula for abstract martingales.

\section*{Supporting Answer}

Elliott's abstract says that the paper constructs an operator whose modulus
does not satisfy the Riesz--Kantorovich formula, thereby answering the
long-standing question.  On the second journal page, the introduction states
the main result explicitly:
\begin{quote}
There exist a compact Hausdorff space $K$ and a regular linear operator
$R:L^1[0,1]\to C(K)$ having a modulus which does not satisfy the
Riesz--Kantorovich formula.
\end{quote}
The same introduction identifies the long-standing question, cites
Aliprantis--Tourky's monograph, and says that the theorem resolves it.  Hence
the general question recalled in arXiv:1501.01685 has a negative answer.

\section*{Scope}

This identification settles precisely the general ``whenever the modulus
exists'' question.  Elliott separately notes that his operator space is not a
vector lattice and leaves open the narrower structural variant asking whether
the formula must hold when the full space of regular operators is a vector
lattice.  This packet makes no claim about that variant.  The
martingale-specific open questions advertised in arXiv:1501.01685 are answered
within that same paper and therefore are extraction false positives rather
than separate literature results.

\section*{Search Evidence}

Cheap run indexes were searched for the source arXiv id, title, regular
abstract martingales, Riesz--Kantorovich, and modulus of an operator; no
duplicate packet was found.  A bounded search for the exact question and
counterexamples located Elliott's publisher record and decisive article.

\begin{thebibliography}{9}

\bibitem{TX15}
V.~G. Troitsky and F.~Xanthos,
\emph{Spaces of regular abstract martingales},
arXiv:1501.01685 (2015).

\bibitem{Elliott19}
M.~Elliott,
\emph{The Riesz--Kantorovich formulae},
Positivity \textbf{23} (2019), 1245--1259,
\href{https://doi.org/10.1007/s11117-019-00661-9}{doi:10.1007/s11117-019-00661-9}.

\end{thebibliography}

\end{document}

